% Options for packages loaded elsewhere
\PassOptionsToPackage{unicode}{hyperref}
\PassOptionsToPackage{hyphens}{url}
\documentclass[
]{article}
\usepackage{xcolor}
\usepackage{amsmath,amssymb}
\setcounter{secnumdepth}{-\maxdimen} % remove section numbering
\usepackage{iftex}
\ifPDFTeX
  \usepackage[T1]{fontenc}
  \usepackage[utf8]{inputenc}
  \usepackage{textcomp} % provide euro and other symbols
\else % if luatex or xetex
  \usepackage{unicode-math} % this also loads fontspec
  \defaultfontfeatures{Scale=MatchLowercase}
  \defaultfontfeatures[\rmfamily]{Ligatures=TeX,Scale=1}
\fi
\usepackage{lmodern}
\ifPDFTeX\else
  % xetex/luatex font selection
\fi
% Use upquote if available, for straight quotes in verbatim environments
\IfFileExists{upquote.sty}{\usepackage{upquote}}{}
\IfFileExists{microtype.sty}{% use microtype if available
  \usepackage[]{microtype}
  \UseMicrotypeSet[protrusion]{basicmath} % disable protrusion for tt fonts
}{}
\makeatletter
\@ifundefined{KOMAClassName}{% if non-KOMA class
  \IfFileExists{parskip.sty}{%
    \usepackage{parskip}
  }{% else
    \setlength{\parindent}{0pt}
    \setlength{\parskip}{6pt plus 2pt minus 1pt}}
}{% if KOMA class
  \KOMAoptions{parskip=half}}
\makeatother
\usepackage{color}
\usepackage{fancyvrb}
\newcommand{\VerbBar}{|}
\newcommand{\VERB}{\Verb[commandchars=\\\{\}]}
\DefineVerbatimEnvironment{Highlighting}{Verbatim}{commandchars=\\\{\}}
% Add ',fontsize=\small' for more characters per line
\newenvironment{Shaded}{}{}
\newcommand{\AlertTok}[1]{\textcolor[rgb]{1.00,0.00,0.00}{\textbf{#1}}}
\newcommand{\AnnotationTok}[1]{\textcolor[rgb]{0.38,0.63,0.69}{\textbf{\textit{#1}}}}
\newcommand{\AttributeTok}[1]{\textcolor[rgb]{0.49,0.56,0.16}{#1}}
\newcommand{\BaseNTok}[1]{\textcolor[rgb]{0.25,0.63,0.44}{#1}}
\newcommand{\BuiltInTok}[1]{\textcolor[rgb]{0.00,0.50,0.00}{#1}}
\newcommand{\CharTok}[1]{\textcolor[rgb]{0.25,0.44,0.63}{#1}}
\newcommand{\CommentTok}[1]{\textcolor[rgb]{0.38,0.63,0.69}{\textit{#1}}}
\newcommand{\CommentVarTok}[1]{\textcolor[rgb]{0.38,0.63,0.69}{\textbf{\textit{#1}}}}
\newcommand{\ConstantTok}[1]{\textcolor[rgb]{0.53,0.00,0.00}{#1}}
\newcommand{\ControlFlowTok}[1]{\textcolor[rgb]{0.00,0.44,0.13}{\textbf{#1}}}
\newcommand{\DataTypeTok}[1]{\textcolor[rgb]{0.56,0.13,0.00}{#1}}
\newcommand{\DecValTok}[1]{\textcolor[rgb]{0.25,0.63,0.44}{#1}}
\newcommand{\DocumentationTok}[1]{\textcolor[rgb]{0.73,0.13,0.13}{\textit{#1}}}
\newcommand{\ErrorTok}[1]{\textcolor[rgb]{1.00,0.00,0.00}{\textbf{#1}}}
\newcommand{\ExtensionTok}[1]{#1}
\newcommand{\FloatTok}[1]{\textcolor[rgb]{0.25,0.63,0.44}{#1}}
\newcommand{\FunctionTok}[1]{\textcolor[rgb]{0.02,0.16,0.49}{#1}}
\newcommand{\ImportTok}[1]{\textcolor[rgb]{0.00,0.50,0.00}{\textbf{#1}}}
\newcommand{\InformationTok}[1]{\textcolor[rgb]{0.38,0.63,0.69}{\textbf{\textit{#1}}}}
\newcommand{\KeywordTok}[1]{\textcolor[rgb]{0.00,0.44,0.13}{\textbf{#1}}}
\newcommand{\NormalTok}[1]{#1}
\newcommand{\OperatorTok}[1]{\textcolor[rgb]{0.40,0.40,0.40}{#1}}
\newcommand{\OtherTok}[1]{\textcolor[rgb]{0.00,0.44,0.13}{#1}}
\newcommand{\PreprocessorTok}[1]{\textcolor[rgb]{0.74,0.48,0.00}{#1}}
\newcommand{\RegionMarkerTok}[1]{#1}
\newcommand{\SpecialCharTok}[1]{\textcolor[rgb]{0.25,0.44,0.63}{#1}}
\newcommand{\SpecialStringTok}[1]{\textcolor[rgb]{0.73,0.40,0.53}{#1}}
\newcommand{\StringTok}[1]{\textcolor[rgb]{0.25,0.44,0.63}{#1}}
\newcommand{\VariableTok}[1]{\textcolor[rgb]{0.10,0.09,0.49}{#1}}
\newcommand{\VerbatimStringTok}[1]{\textcolor[rgb]{0.25,0.44,0.63}{#1}}
\newcommand{\WarningTok}[1]{\textcolor[rgb]{0.38,0.63,0.69}{\textbf{\textit{#1}}}}
\usepackage{longtable,booktabs,array}
\newcounter{none} % for unnumbered tables
\usepackage{calc} % for calculating minipage widths
% Correct order of tables after \paragraph or \subparagraph
\usepackage{etoolbox}
\makeatletter
\patchcmd\longtable{\par}{\if@noskipsec\mbox{}\fi\par}{}{}
\makeatother
% Allow footnotes in longtable head/foot
\IfFileExists{footnotehyper.sty}{\usepackage{footnotehyper}}{\usepackage{footnote}}
\makesavenoteenv{longtable}
\usepackage{graphicx}
\makeatletter
\newsavebox\pandoc@box
\newcommand*\pandocbounded[1]{% scales image to fit in text height/width
  \sbox\pandoc@box{#1}%
  \Gscale@div\@tempa{\textheight}{\dimexpr\ht\pandoc@box+\dp\pandoc@box\relax}%
  \Gscale@div\@tempb{\linewidth}{\wd\pandoc@box}%
  \ifdim\@tempb\p@<\@tempa\p@\let\@tempa\@tempb\fi% select the smaller of both
  \ifdim\@tempa\p@<\p@\scalebox{\@tempa}{\usebox\pandoc@box}%
  \else\usebox{\pandoc@box}%
  \fi%
}
% Set default figure placement to htbp
\def\fps@figure{htbp}
\makeatother
\setlength{\emergencystretch}{3em} % prevent overfull lines
\providecommand{\tightlist}{%
  \setlength{\itemsep}{0pt}\setlength{\parskip}{0pt}}
\usepackage{bookmark}
\IfFileExists{xurl.sty}{\usepackage{xurl}}{} % add URL line breaks if available
\urlstyle{same}
\hypersetup{
  hidelinks,
  pdfcreator={LaTeX via pandoc}}

\author{}
\date{}

\begin{document}

\section{Processing Appunti 2 e 3 videolezione di
affiancamento.}\label{processing-appunti-2-e-3-videolezione-di-affiancamento}

\begin{Shaded}
\begin{Highlighting}[]
\FunctionTok{point}\OperatorTok{();} \CommentTok{//per fare un punto}
\FunctionTok{line}\OperatorTok{();} \CommentTok{//per fare una linea}
\FunctionTok{ellipse}\OperatorTok{();} \CommentTok{//per fare un ellisse}
\FunctionTok{circle}\OperatorTok{();} \CommentTok{//per fare un cerchio, prende x,y del centro e ed il raggio {-}\textgreater{} 3 parametri}
\FunctionTok{rect}\OperatorTok{();} \CommentTok{//per fare u rettangolo}
\FunctionTok{quad}\OperatorTok{();} \CommentTok{//per fare un quadrilatero (4 vertici)}
\FunctionTok{shape}\OperatorTok{();} \CommentTok{//per fare una forma con i vertici}
\end{Highlighting}
\end{Shaded}

\begin{Shaded}
\begin{Highlighting}[]
\CommentTok{//[PRIMO FILE]}
\FunctionTok{size}\OperatorTok{(}\DecValTok{500}\OperatorTok{,}\DecValTok{500}\OperatorTok{);}
\CommentTok{//L\textquotesingle{}arco verrà disegnato in senso orario, quindi parte da un punto e torna all\textquotesingle{}altrosize(500,500);}
\CommentTok{//Disegno dimostrativo di tutti e 4 gli spicchi, di ogni metà, e della circonferenza chiusa}
\CommentTok{//arc(a,b,c,d,start,stop) {-}\textgreater{}ATTENZIONE: Stop \textgreater{} Start!!! deve essere così.}
\CommentTok{//a e b sono le coordinate del centro, c e d sono le dimensioni di larghezza e altezza, e start e stop sono angoli (in radianti) di inzio e fine.}

\CommentTok{/* In Processing, la funzione arc() ha un comportamento rigido:}
\CommentTok{{-} Il punto 0 (zero radianti) è a ore 3:00 (a destra).}
\CommentTok{{-} Disegna sempre in senso ORARIO (da Start verso Stop).}
\CommentTok{{-} Disegna solo se Stop \textgreater{} Start. */}

\FunctionTok{arc}\OperatorTok{(}\DecValTok{60}\OperatorTok{,} \DecValTok{60}\OperatorTok{,} \DecValTok{80}\OperatorTok{,} \DecValTok{80}\OperatorTok{,}\NormalTok{ PI}\OperatorTok{,}\NormalTok{ PI }\OperatorTok{+}\NormalTok{ HALF\_PI}\OperatorTok{);} \CommentTok{//1 quadrante}
\FunctionTok{arc}\OperatorTok{(}\DecValTok{160}\OperatorTok{,} \DecValTok{60}\OperatorTok{,} \DecValTok{80}\OperatorTok{,} \DecValTok{80}\OperatorTok{,}\NormalTok{ PI }\OperatorTok{+}\NormalTok{ HALF\_PI}\OperatorTok{,}\NormalTok{ TWO\_PI}\OperatorTok{);} \CommentTok{//2 quadrante}
\FunctionTok{arc}\OperatorTok{(}\DecValTok{260}\OperatorTok{,} \DecValTok{60}\OperatorTok{,} \DecValTok{80}\OperatorTok{,} \DecValTok{80}\OperatorTok{,} \DecValTok{0}\OperatorTok{,}\NormalTok{ HALF\_PI}\OperatorTok{);} \CommentTok{//3 quadrane}
\FunctionTok{arc}\OperatorTok{(}\DecValTok{360}\OperatorTok{,} \DecValTok{60}\OperatorTok{,} \DecValTok{80}\OperatorTok{,} \DecValTok{80}\OperatorTok{,}\NormalTok{ HALF\_PI}\OperatorTok{,}\NormalTok{ PI}\OperatorTok{);} \CommentTok{//4 quadrante}
\CommentTok{/*Il punto "Ore 12:00" ha due nomi:}
\CommentTok{Nome lungo (giro completo): 270° (PI + HALF\_PI).}
\CommentTok{Nome corto (indietro): {-}90° ({-}HALF\_PI).*/}


\FunctionTok{arc}\OperatorTok{(}\DecValTok{60}\OperatorTok{,} \DecValTok{200}\OperatorTok{,} \DecValTok{80}\OperatorTok{,} \DecValTok{80}\OperatorTok{,}\NormalTok{ PI}\OperatorTok{,}\NormalTok{ TWO\_PI}\OperatorTok{);} \CommentTok{//metà in alto}
\FunctionTok{arc}\OperatorTok{(}\DecValTok{160}\OperatorTok{,} \DecValTok{200}\OperatorTok{,} \DecValTok{80}\OperatorTok{,} \DecValTok{80}\OperatorTok{,} \DecValTok{0}\OperatorTok{,}\NormalTok{ PI}\OperatorTok{);} \CommentTok{//metà in basso}
\FunctionTok{arc}\OperatorTok{(}\DecValTok{260}\OperatorTok{,} \DecValTok{200}\OperatorTok{,} \DecValTok{80}\OperatorTok{,} \DecValTok{80}\OperatorTok{,}\NormalTok{ HALF\_PI}\OperatorTok{,}\NormalTok{ PI }\OperatorTok{+}\NormalTok{ HALF\_PI}\OperatorTok{);} \CommentTok{//metà a sinistra}
\FunctionTok{arc}\OperatorTok{(}\DecValTok{360}\OperatorTok{,} \DecValTok{200}\OperatorTok{,} \DecValTok{80}\OperatorTok{,} \DecValTok{80}\OperatorTok{,} \OperatorTok{{-}}\NormalTok{HALF\_PI}\OperatorTok{,}\NormalTok{ HALF\_PI}\OperatorTok{);} \CommentTok{//metà a destra}

\FunctionTok{arc}\OperatorTok{(}\DecValTok{60}\OperatorTok{,} \DecValTok{340}\OperatorTok{,} \DecValTok{80}\OperatorTok{,} \DecValTok{80}\OperatorTok{,} \DecValTok{0}\OperatorTok{,}\NormalTok{ TWO\_PI}\OperatorTok{);}\NormalTok{\_PI)}\OperatorTok{;}


\CommentTok{//Introduciamo un\textquotesingle{}altra funzione, che ci servirà per disegnare pacman}
\CommentTok{//ossia la funzione \textquotesingle{}radins()\textquotesingle{}, se inseriamo un angolo in gradi lo converte in radianti}
\CommentTok{//quindi pacman sarebbe:}
\FunctionTok{arc}\OperatorTok{(}\NormalTok{width}\OperatorTok{/}\DecValTok{2}\OperatorTok{,}\NormalTok{ height}\OperatorTok{/}\DecValTok{2}\OperatorTok{,} \DecValTok{100}\OperatorTok{,} \DecValTok{100}\OperatorTok{,} \FunctionTok{radians}\OperatorTok{(}\DecValTok{30}\OperatorTok{),} \FunctionTok{radians}\OperatorTok{(}\DecValTok{330}\OperatorTok{));}


\CommentTok{//Però la nostra "torta" rimarrà aperta e senza stroke che la chiude, per questo \textquotesingle{}arc()\textquotesingle{} accetta un settimo parameto che sarebbe la MODE, e ne abbiamo 3}
\CommentTok{/*}
\CommentTok{{-} PIE: chiude la nostra "torta/pie"}
\CommentTok{{-} OPEN: lascia la linea aperta, ma non ha la forma a torta, collega i due punti con una retta, ma rimane aperta}
\CommentTok{{-} CHORD: collega i due punti con una linea chiusa, ossia con un suo stroke}
\CommentTok{Dunque il nostro pacman sarebbe:}
\CommentTok{*/}

\FunctionTok{arc}\OperatorTok{(}\NormalTok{width}\OperatorTok{/}\DecValTok{2}\OperatorTok{,}\NormalTok{ height}\OperatorTok{/}\DecValTok{2}\OperatorTok{,} \DecValTok{100}\OperatorTok{,} \DecValTok{100}\OperatorTok{,} \FunctionTok{radians}\OperatorTok{(}\DecValTok{30}\OperatorTok{),} \FunctionTok{radians}\OperatorTok{(}\DecValTok{330}\OperatorTok{),}\NormalTok{ PIE}\OperatorTok{);}
\end{Highlighting}
\end{Shaded}

\pandocbounded{\includegraphics[keepaspectratio,alt={}]{C:/Users/giuli/AppData/Roaming/marktext/images/2026-02-12-02-08-08-image.png}}

\begin{Shaded}
\begin{Highlighting}[]
\CommentTok{//[SECONDO FILE]}
\CommentTok{//Vediamo l\textquotesingle{}ultima funzione di disegno che ci rimane}
\CommentTok{//E la funzione è: \textquotesingle{}bezier()\textquotesingle{}}

\CommentTok{/*Disegna una curva di Bézier sullo schermo.}
\CommentTok{Queste curve sono definite da una serie di punti di ancoraggio e di controllo.}
\CommentTok{I primi due parametri specificano il primo punto di ancoraggio e gli ultimi due parametri specificano l\textquotesingle{}altro punto di ancoraggio.}
\CommentTok{I parametri centrali specificano i punti di controllo che definiscono la forma della curva. */}

\DataTypeTok{void} \FunctionTok{setup}\OperatorTok{()\{}
  \FunctionTok{size}\OperatorTok{(}\DecValTok{500}\OperatorTok{,}\DecValTok{500}\OperatorTok{);}
  \FunctionTok{noFill}\OperatorTok{();}
\OperatorTok{\}}

\DataTypeTok{void} \FunctionTok{draw}\OperatorTok{()\{}
  \FunctionTok{background}\OperatorTok{(}\DecValTok{255}\OperatorTok{);} \CommentTok{//Reimpostare il background serve per rendere il disegno una vera e propria animazione e non un disegno ripetuto uno sopra l\textquotesingle{}altro. Prova a rimuover questa istruzione per capire meglio.}
  \FunctionTok{stroke}\OperatorTok{(}\DecValTok{255}\OperatorTok{,}\DecValTok{100}\OperatorTok{,}\DecValTok{0}\OperatorTok{);}
  \FunctionTok{line}\OperatorTok{(}\DecValTok{185}\OperatorTok{,}\DecValTok{50}\OperatorTok{,}\NormalTok{ mouseX}\OperatorTok{,}\NormalTok{ mouseY}\OperatorTok{);}
  \FunctionTok{line}\OperatorTok{(}\DecValTok{15}\OperatorTok{,}\DecValTok{290}\OperatorTok{,} \DecValTok{190}\OperatorTok{,}\DecValTok{300}\OperatorTok{);}
  \FunctionTok{stroke}\OperatorTok{(}\DecValTok{0}\OperatorTok{);}
  \CommentTok{//bezier(ancorPoint1\_X, ancorPoint2\_Y, controlPoint1\_X, controlPoint1\_Y, controlPoint2\_X, controlPoint2\_Y, ancorPoint1\_X, ancorPoint1\_Y);}
  \FunctionTok{bezier}\OperatorTok{(}\DecValTok{185}\OperatorTok{,}\DecValTok{50}\OperatorTok{,}\NormalTok{ mouseX}\OperatorTok{,}\NormalTok{mouseY}\OperatorTok{,} \DecValTok{190}\OperatorTok{,}\DecValTok{300}\OperatorTok{,} \DecValTok{15}\OperatorTok{,}\DecValTok{290}\OperatorTok{);}

\OperatorTok{\}}
\end{Highlighting}
\end{Shaded}

\begin{Shaded}
\begin{Highlighting}[]
\CommentTok{//[TERZO FILE]}

\CommentTok{//Facciamo la prima animazione}
\DataTypeTok{int}\NormalTok{ pos1 }\OperatorTok{=} \DecValTok{0}\OperatorTok{;}
\DataTypeTok{int}\NormalTok{ pos2 }\OperatorTok{=} \DecValTok{0}\OperatorTok{;}

\DataTypeTok{void} \FunctionTok{setup}\OperatorTok{()\{}
  \FunctionTok{size}\OperatorTok{(}\DecValTok{500}\OperatorTok{,}\DecValTok{500}\OperatorTok{);}

\OperatorTok{\}}


\DataTypeTok{void} \FunctionTok{draw}\OperatorTok{()\{}
  \FunctionTok{background}\OperatorTok{(}\DecValTok{255}\OperatorTok{);}
  \FunctionTok{line}\OperatorTok{(}\NormalTok{pos1}\OperatorTok{,}\DecValTok{20}\OperatorTok{,}\NormalTok{ pos2}\OperatorTok{,}\NormalTok{height}\OperatorTok{);}
\NormalTok{  pos1}\OperatorTok{++;}
\NormalTok{  pos2}\OperatorTok{++;}
\OperatorTok{\}}

\CommentTok{//Animazione di una semplice linea che si sposta da sinistra a destra}
\end{Highlighting}
\end{Shaded}

\begin{Shaded}
\begin{Highlighting}[]
\CommentTok{//[QUARTO FILE]}

\DataTypeTok{int}\NormalTok{ pos }\OperatorTok{=} \DecValTok{0}\OperatorTok{;} \CommentTok{//Come sempre rendiamo globale lo scope della variabile che ci serve usare}

\DataTypeTok{void} \FunctionTok{setup}\OperatorTok{()\{}
  \FunctionTok{size}\OperatorTok{(}\DecValTok{500}\OperatorTok{,}\DecValTok{500}\OperatorTok{);}
  \FunctionTok{frameRate}\OperatorTok{(}\DecValTok{200}\OperatorTok{);}
\OperatorTok{\}}

\CommentTok{/*}
\CommentTok{In questo esercizio vogliamo una linea che si muove da sinistra a destra facendo l\textquotesingle{}effetto pacman, ossia che una volta arrivato alla fine dello schermo appare di nuovo a sinistra all\textquotesingle{}infinito.}
\CommentTok{Ovviamente impostiamo il nostro background ad ogni iterazione della funzione drawInizializziamo la nostra linea, una semplice linea verticale che attraversa tutta la canvas.}
\CommentTok{Adesso presta attenzione.}
\CommentTok{Se il nostro incremento sarà pos++ (quindi che aumenta la coordinata x di 1) allora il funzionerà tutto}
\CommentTok{Ma se il nostro incremento sarà pos+=3 allora non passeremo mai a pos == 500 (che è la nostra width) e dunque la condizione non sarà mai vera. Per questo è opportuno modificare la condizione come pos \textgreater{}= width, ciò vuol dire che se pos sarà comunque uguale o poco maggiore di 500 dovrà ripartire comunque da 0.}
\CommentTok{*/}

\CommentTok{//SCRIVO DUE FUNZIONI DRAW PER MOSTRARE LE DIFFERENZE, MA OVVIAMENTE POSSIAMO AVERNE SOLAMENTE UNA!!!}

\DataTypeTok{void} \FunctionTok{draw}\OperatorTok{()\{}
  \FunctionTok{background}\OperatorTok{(}\DecValTok{255}\OperatorTok{);}
  \FunctionTok{line}\OperatorTok{(}\NormalTok{pos}\OperatorTok{,}\DecValTok{0}\OperatorTok{,}\NormalTok{ pos}\OperatorTok{,}\NormalTok{height}\OperatorTok{);}
\NormalTok{  pos}\OperatorTok{++;}
  \ControlFlowTok{if}\OperatorTok{(}\NormalTok{pos }\OperatorTok{==}\NormalTok{ width}\OperatorTok{)\{}
\NormalTok{    pos }\OperatorTok{=} \DecValTok{0}\OperatorTok{;}
  \OperatorTok{\}}
\OperatorTok{\}}

\DataTypeTok{void} \FunctionTok{draw}\OperatorTok{()\{}
  \FunctionTok{background}\OperatorTok{(}\DecValTok{255}\OperatorTok{);}
  \FunctionTok{line}\OperatorTok{(}\NormalTok{pos}\OperatorTok{,}\DecValTok{0}\OperatorTok{,}\NormalTok{ pos}\OperatorTok{,}\NormalTok{height}\OperatorTok{);}
\NormalTok{  pos}\OperatorTok{+=}\DecValTok{3}\OperatorTok{;}
  \ControlFlowTok{if}\OperatorTok{(}\NormalTok{pos }\OperatorTok{\textgreater{}=}\NormalTok{ width}\OperatorTok{)\{}
\NormalTok{    pos }\OperatorTok{=} \DecValTok{0}\OperatorTok{;}
  \OperatorTok{\}}
\OperatorTok{\}}
\end{Highlighting}
\end{Shaded}

\begin{Shaded}
\begin{Highlighting}[]
\CommentTok{//[QUINTO FILE]}
\CommentTok{//Adesso vogliamo eliminare l\textquotesingle{}effetto pacman, e fare in modo che la linea rimbalzi a destra e sinista}

\DataTypeTok{int}\NormalTok{ pos }\OperatorTok{=} \DecValTok{0}\OperatorTok{;}
\DataTypeTok{int}\NormalTok{ dir }\OperatorTok{=} \DecValTok{1}\OperatorTok{;} \CommentTok{//direzione}
\DataTypeTok{int}\NormalTok{ v }\OperatorTok{=} \DecValTok{3}\OperatorTok{;} \CommentTok{//velocità}

\DataTypeTok{void} \FunctionTok{setup}\OperatorTok{()\{}
  \FunctionTok{size}\OperatorTok{(}\DecValTok{500}\OperatorTok{,}\DecValTok{500}\OperatorTok{);}

\OperatorTok{\}}

\DataTypeTok{void} \FunctionTok{draw}\OperatorTok{()\{}
  \FunctionTok{background}\OperatorTok{(}\DecValTok{255}\OperatorTok{);}
  \FunctionTok{line}\OperatorTok{(}\NormalTok{pos}\OperatorTok{,}\DecValTok{20}\OperatorTok{,}\NormalTok{ pos}\OperatorTok{,}\NormalTok{height}\OperatorTok{);}
  \CommentTok{//pos+=3 vuol dire pos += direzione * velocità, perchè si sposta di 3 per un framerate di 3}
  \CommentTok{//per fare avanti ed indietro dobbiamo solo aumentare e diminuire e poi aumentare e così via la dimensione.}
   \ControlFlowTok{if}\OperatorTok{(}\NormalTok{pos}\OperatorTok{\textgreater{}}\NormalTok{width }\OperatorTok{||}\NormalTok{ pos}\OperatorTok{\textless{}}\DecValTok{0}\OperatorTok{)\{}
\NormalTok{   dir }\OperatorTok{=} \OperatorTok{{-}}\NormalTok{dir}\OperatorTok{;}
   \OperatorTok{\}}

\NormalTok{   pos }\OperatorTok{+=}\NormalTok{ dir }\OperatorTok{*}\NormalTok{ v}\OperatorTok{;}
\OperatorTok{\}}


\CommentTok{/*1. Il Segreto del segno Meno ({-})}
\CommentTok{In matematica e informatica, moltiplicare per {-}1 è l\textquotesingle{}equivalente di fare una inversione a U con l\textquotesingle{}auto.}
\CommentTok{Guarda cosa succede ai numeri se li moltiplichi per {-}1:}
\CommentTok{Se hai 1 (stai andando avanti) e lo moltiplichi per {-}1, ottieni {-}1 (vai indietro).}
\CommentTok{Se hai {-}1 (stai andando indietro) e lo moltiplichi per {-}1, ottieni 1 (meno per meno fa più {-}\textgreater{} torni ad andare avanti).}

\CommentTok{2. Cosa succede nel tuo codice?}
\CommentTok{Tu hai questa riga fondamentale alla fine: pos += dir * v;}
\CommentTok{Analizziamola in due scenari:}
\CommentTok{SCENARIO A: La marcia avanti (dir = 1) Stai iniziando il programma. dir è 1. Il calcolo diventa: pos = pos + (1 * 3). Ovvero: pos = pos + 3. La linea si sposta a destra: 0, 3, 6, 9... fino a 500.}
\CommentTok{SCENARIO B: Il Muro (dir = {-}dir) Appena tocchi il bordo (pos \textgreater{} width), il codice esegue: dir = {-}dir; Siccome dir era 1, ora diventa {-}1.}
\CommentTok{SCENARIO C: La retromarcia (dir = {-}1) Al frame successivo, il codice riesegue il calcolo di movimento: pos += dir * v; Sostituiamo i numeri: pos = pos + ({-}1 * 3). Ovvero: pos = pos {-} 3. Vedi? Ora sottrai 3! La linea torna indietro: 500, 497, 494...}

\CommentTok{In sintesi}
\CommentTok{Quella riga dir = {-}dir è un interruttore. È un modo velocissimo per dire al computer: "Qualunque cosa tu sia (1 o {-}1), diventa il tuo opposto".}
\CommentTok{Se sei positivo, diventa negativo.}
\CommentTok{Se sei negativo, diventa positivo.}
\CommentTok{È come sbattere contro un muro e rimbalzare: la tua velocità (v) è sempre 3 km/h, ma la tua direzione (dir) cambia verso.*/}


\CommentTok{//Rifacciamo la stessa animazione ma con la riga che si muove in orizzontale}
\DataTypeTok{int}\NormalTok{ pos }\OperatorTok{=} \DecValTok{0}\OperatorTok{;}
\DataTypeTok{int}\NormalTok{ dir }\OperatorTok{=} \DecValTok{1}\OperatorTok{;}
\DataTypeTok{int}\NormalTok{ v }\OperatorTok{=} \DecValTok{3}\OperatorTok{;}

\DataTypeTok{void} \FunctionTok{setup}\OperatorTok{()\{}
  \FunctionTok{size}\OperatorTok{(}\DecValTok{500}\OperatorTok{,} \DecValTok{500}\OperatorTok{);}
\OperatorTok{\}}

\DataTypeTok{void} \FunctionTok{draw}\OperatorTok{()\{}
  \FunctionTok{background}\OperatorTok{(}\DecValTok{255}\OperatorTok{);}
  \FunctionTok{strokeWeight}\OperatorTok{(}\DecValTok{2}\OperatorTok{);}
  \FunctionTok{line}\OperatorTok{(}\DecValTok{0}\OperatorTok{,}\NormalTok{pos}\OperatorTok{,}\NormalTok{ width}\OperatorTok{,}\NormalTok{pos}\OperatorTok{);}

  \ControlFlowTok{if}\OperatorTok{(}\NormalTok{pos}\OperatorTok{\textgreater{}}\NormalTok{height }\OperatorTok{||}\NormalTok{ pos}\OperatorTok{\textless{}}\DecValTok{0}\OperatorTok{)\{} \CommentTok{//Se il valore di y è \textgreater{} 500 o \textless{} di 0, allora rendi la direzione negativa {-}\textgreater{} torna indietro}
\NormalTok{    dir }\OperatorTok{=} \OperatorTok{{-}}\NormalTok{dir}\OperatorTok{;}
  \OperatorTok{\}}

\NormalTok{  pos }\OperatorTok{+=}\NormalTok{ dir}\OperatorTok{*}\NormalTok{v}\OperatorTok{;} \CommentTok{//per far muovere la linea si fa la posizione = posizione + (dir * v)}
  \CommentTok{//Ad esempio {-}\textgreater{} posizione (2) = 2 + (+1 * 3); }
\OperatorTok{\}}
\end{Highlighting}
\end{Shaded}

\begin{Shaded}
\begin{Highlighting}[]
\CommentTok{//[SESTO FILE]}
\CommentTok{//Esercizio di mattonelle sconnesse}
\FunctionTok{size}\OperatorTok{(}\DecValTok{500}\OperatorTok{,}\DecValTok{500}\OperatorTok{);}
\FunctionTok{background}\OperatorTok{(}\DecValTok{0}\OperatorTok{,}\DecValTok{255}\OperatorTok{,}\DecValTok{0}\OperatorTok{);}
\FunctionTok{noStroke}\OperatorTok{();}

\CommentTok{//Usiamo la funzione \textquotesingle{}random()\textquotesingle{}, prende uno o due parametri.}
\CommentTok{//Se usiamo random(30) ci restituisce un float. Quindi da 0.0 a 29.999... Quindi non avremo mai 30}
\CommentTok{//Il casting in float si fa come nel C: (int)random(30)}
\CommentTok{//Se passiamo due parametri stiamo attivando un range più definito. Quindi se usiamo:}
\CommentTok{//random(10,15) avremo il range che va da 10.0 a 14.999...}

\DataTypeTok{int}\NormalTok{ l}\OperatorTok{=}\DecValTok{50}\OperatorTok{;} \CommentTok{//lato della mattonella}

\ControlFlowTok{for}\OperatorTok{(}\DataTypeTok{int}\NormalTok{ y}\OperatorTok{=}\DecValTok{0}\OperatorTok{;}\NormalTok{ y}\OperatorTok{\textless{}}\NormalTok{height}\OperatorTok{;}\NormalTok{ y}\OperatorTok{+=}\NormalTok{l}\OperatorTok{)\{}
  \ControlFlowTok{for}\OperatorTok{(}\DataTypeTok{int}\NormalTok{ x}\OperatorTok{=}\DecValTok{0}\OperatorTok{;}\NormalTok{ x}\OperatorTok{\textless{}}\NormalTok{width}\OperatorTok{;}\NormalTok{x}\OperatorTok{+=}\NormalTok{l}\OperatorTok{)\{}
    \FunctionTok{fill}\OperatorTok{(}\FunctionTok{random}\OperatorTok{(}\DecValTok{255}\OperatorTok{),}\DecValTok{130}\OperatorTok{);} \CommentTok{//Questa riga riempie casualmente e con varie tonalità i colori. Perché impostiamo livelli di grigio diversi di trasparenze diverse.}
    \FunctionTok{rect}\OperatorTok{(}\NormalTok{x}\OperatorTok{+}\FunctionTok{random}\OperatorTok{({-}}\DecValTok{5}\OperatorTok{,}\DecValTok{5}\OperatorTok{),}\NormalTok{y}\OperatorTok{+}\FunctionTok{random}\OperatorTok{({-}}\DecValTok{5}\OperatorTok{,+}\DecValTok{5}\OperatorTok{),}\NormalTok{l}\OperatorTok{,}\NormalTok{l}\OperatorTok{);} \CommentTok{//come al solito, i primi due parametri impostano l\textquotesingle{}angolo in alto a sinistra e gli altri due larghezza e altezza}
    \CommentTok{//Aggiungiamo \textquotesingle{}x+random({-}5,5),y+random({-}5,+5)\textquotesingle{}. Quello che abbiamo fatto è stato spostare di 5 pixel a sinistra o destra o in alto o in basso}
  \OperatorTok{\}}
\OperatorTok{\}}

\CommentTok{//Ricordiamo i parametri accettati da \textquotesingle{}fill()\textquotesingle{}}
\CommentTok{/*}
\CommentTok{{-} 3 parametri = RGB}
\CommentTok{{-} 1 paramrtro = livello di grigio}
\CommentTok{{-} 1 parametro = \#esadeciamle}
\CommentTok{{-} 4 parametri = RGB e alpga}
\CommentTok{{-} 2 parametri = livello di grigio e alpha}
\CommentTok{*/}
\end{Highlighting}
\end{Shaded}

\begin{Shaded}
\begin{Highlighting}[]
\CommentTok{//[SETTIMO FILE]}
\DataTypeTok{void} \FunctionTok{setup}\OperatorTok{()\{}
\FunctionTok{size}\OperatorTok{(}\DecValTok{500}\OperatorTok{,}\DecValTok{500}\OperatorTok{);}
\FunctionTok{background}\OperatorTok{(}\DecValTok{0}\OperatorTok{,}\DecValTok{255}\OperatorTok{,}\DecValTok{0}\OperatorTok{);}
\DataTypeTok{int}\NormalTok{ l}\OperatorTok{=}\DecValTok{50}\OperatorTok{;}
\ControlFlowTok{for}\OperatorTok{(}\DataTypeTok{int}\NormalTok{ y}\OperatorTok{=}\DecValTok{0}\OperatorTok{;}\NormalTok{ y}\OperatorTok{\textless{}}\NormalTok{height}\OperatorTok{;}\NormalTok{ y}\OperatorTok{+=}\NormalTok{l}\OperatorTok{)\{}
  \ControlFlowTok{for}\OperatorTok{(}\DataTypeTok{int}\NormalTok{ x}\OperatorTok{=}\DecValTok{0}\OperatorTok{;}\NormalTok{ x}\OperatorTok{\textless{}}\NormalTok{width}\OperatorTok{;}\NormalTok{x}\OperatorTok{+=}\NormalTok{l}\OperatorTok{)\{}
    \FunctionTok{fill}\OperatorTok{(}\FunctionTok{random}\OperatorTok{(}\DecValTok{255}\OperatorTok{),}\DecValTok{130}\OperatorTok{);} 
    \FunctionTok{rect}\OperatorTok{(}\NormalTok{x}\OperatorTok{+}\FunctionTok{random}\OperatorTok{({-}}\DecValTok{5}\OperatorTok{,}\DecValTok{5}\OperatorTok{),}\NormalTok{y}\OperatorTok{+}\FunctionTok{random}\OperatorTok{({-}}\DecValTok{5}\OperatorTok{,+}\DecValTok{5}\OperatorTok{),}\NormalTok{l}\OperatorTok{,}\NormalTok{l}\OperatorTok{);}
    \OperatorTok{\}}
  \OperatorTok{\}}
\OperatorTok{\}}

\DataTypeTok{void} \FunctionTok{draw}\OperatorTok{()\{}
\OperatorTok{\}}

\DataTypeTok{void} \FunctionTok{mousePressed}\OperatorTok{()\{} \CommentTok{//Funzione che capisce se premi un tasto del mouse}
\FunctionTok{setup}\OperatorTok{();}\CommentTok{//ogni volta che premiamo un tasto richiama la funzione setup}
\OperatorTok{\}}
\end{Highlighting}
\end{Shaded}

\begin{Shaded}
\begin{Highlighting}[]
\CommentTok{//[OTTAVO FILE]}
\DataTypeTok{void} \FunctionTok{setup}\OperatorTok{()\{}
  \FunctionTok{size}\OperatorTok{(}\DecValTok{500}\OperatorTok{,}\DecValTok{500}\OperatorTok{);}
\OperatorTok{\}}

\DataTypeTok{void} \FunctionTok{draw}\OperatorTok{()\{}
  \ControlFlowTok{if}\OperatorTok{(}\NormalTok{mousePressed}\OperatorTok{)\{} \CommentTok{//E\textquotesingle{} una variabile booleana. False se non premiamo nulla, True se premiamo qualcosa.}
    \CommentTok{//mouseButton conserva quale abbiamo premuto}
    \CommentTok{//mouseButton è una variabile di tipo intero}
    \ControlFlowTok{if}\OperatorTok{(}\NormalTok{mouseButton}\OperatorTok{==}\NormalTok{LEFT}\OperatorTok{)\{} \CommentTok{//Se premiamo tasto sinistro stampa L}
      \FunctionTok{println}\OperatorTok{(}\StringTok{"L"}\OperatorTok{);}
  \OperatorTok{\}}
    \ControlFlowTok{if}\OperatorTok{(}\NormalTok{mouseButton}\OperatorTok{==}\NormalTok{RIGHT}\OperatorTok{)\{} \CommentTok{//Se premiamo tasto destro stampa R}
      \FunctionTok{println}\OperatorTok{(}\StringTok{"R"}\OperatorTok{);}
    \OperatorTok{\}}
  \OperatorTok{\}}

  \ControlFlowTok{if}\OperatorTok{(}\NormalTok{keyPressed}\OperatorTok{)\{} \CommentTok{//keyPressed è booleana come mousePressed}
  \CommentTok{//key è una funzione di tipo char}
    \ControlFlowTok{if}\OperatorTok{(}\NormalTok{key }\OperatorTok{==} \CharTok{\textquotesingle{}a\textquotesingle{}} \OperatorTok{||}\NormalTok{ key }\OperatorTok{==} \CharTok{\textquotesingle{}A\textquotesingle{}}\OperatorTok{)\{}
      \FunctionTok{println}\OperatorTok{(}\StringTok{"ho premuto \textquotesingle{}a\textquotesingle{}"}\OperatorTok{);}
    \OperatorTok{\}}
  \OperatorTok{\}}
\OperatorTok{\}}
\end{Highlighting}
\end{Shaded}

\begin{Shaded}
\begin{Highlighting}[]
\CommentTok{//[NONO FILE]}
\CommentTok{//creiamo un array di tipo float globalmente in modo che è accessibile ovunque}
\CommentTok{//new è l\textquotesingle{}equivalente del malloc in C, ed in questo caso abbiamo allocato un array di 10 valori, ossia 10 cerchi}
\DataTypeTok{float} \OperatorTok{[]}\NormalTok{ d }\OperatorTok{=} \KeywordTok{new} \DataTypeTok{float} \OperatorTok{[}\DecValTok{10}\OperatorTok{];} 

\DataTypeTok{void} \FunctionTok{setup}\OperatorTok{()\{}
  \FunctionTok{size}\OperatorTok{(}\DecValTok{800}\OperatorTok{,} \DecValTok{800}\OperatorTok{);}
  \ControlFlowTok{for} \OperatorTok{(}\DataTypeTok{int}\NormalTok{ i }\OperatorTok{=} \DecValTok{0}\OperatorTok{;}\NormalTok{ i }\OperatorTok{\textless{}}\NormalTok{ d}\OperatorTok{.}\FunctionTok{length}\OperatorTok{;}\NormalTok{ i}\OperatorTok{++)\{} \CommentTok{//riempiamo l\textquotesingle{}array}
\NormalTok{    d}\OperatorTok{[}\NormalTok{i}\OperatorTok{]} \OperatorTok{=} \OperatorTok{(}\NormalTok{i}\OperatorTok{+}\DecValTok{1}\OperatorTok{)*}\DecValTok{60}\OperatorTok{;} \CommentTok{//ogni valore sarà quello precedente + 60}
  \OperatorTok{\}}
\OperatorTok{\}}

\DataTypeTok{void} \FunctionTok{draw}\OperatorTok{()\{}
  \ControlFlowTok{for} \OperatorTok{(}\DataTypeTok{int}\NormalTok{ i }\OperatorTok{=}\NormalTok{ d}\OperatorTok{.}\FunctionTok{length}\OperatorTok{{-}}\DecValTok{1}\OperatorTok{;}\NormalTok{ i}\OperatorTok{\textgreater{}=}\DecValTok{0}\OperatorTok{;}\NormalTok{ i}\OperatorTok{{-}{-})\{} \CommentTok{//Scorriamo l\textquotesingle{}array al contrario altrimenti i cerchi più grandi saranno disegnati su quelli più piccoli }
    \FunctionTok{fill}\OperatorTok{(}\FunctionTok{random}\OperatorTok{(}\DecValTok{255}\OperatorTok{),} \FunctionTok{random}\OperatorTok{(}\DecValTok{255}\OperatorTok{),} \FunctionTok{random}\OperatorTok{(}\DecValTok{255}\OperatorTok{));} \CommentTok{//Riempiamo con colori random}
    \FunctionTok{ellipse}\OperatorTok{(}\NormalTok{width}\OperatorTok{/}\DecValTok{2}\OperatorTok{,}\NormalTok{ height}\OperatorTok{/}\DecValTok{2}\OperatorTok{,}\NormalTok{ d}\OperatorTok{[}\NormalTok{i}\OperatorTok{],}\NormalTok{ d}\OperatorTok{[}\NormalTok{i}\OperatorTok{]);} \CommentTok{//Il centro di ogni ellisse sarà al centro della canvas, e i diametri avranno i valori che andiamo a leggere dall\textquotesingle{}array}
  \OperatorTok{\}}

  \ControlFlowTok{if}\OperatorTok{(}\NormalTok{mousePressed}\OperatorTok{)\{}
    \ControlFlowTok{if}\OperatorTok{(}\NormalTok{mouseButton }\OperatorTok{==}\NormalTok{ LEFT}\OperatorTok{)\{}
        \FunctionTok{println}\OperatorTok{(}\StringTok{"L"}\OperatorTok{);}
        \FunctionTok{noLoop}\OperatorTok{();} \CommentTok{//interrompe l\textquotesingle{}esecuzione della funzione \textquotesingle{}draw()\textquotesingle{} se noi premiamo il tasto sinistro del mouse}
    \OperatorTok{\}}
    \ControlFlowTok{if}\OperatorTok{(}\NormalTok{mouseButton }\OperatorTok{==}\NormalTok{ RIGHT}\OperatorTok{)\{}
        \FunctionTok{println}\OperatorTok{(}\StringTok{"R"}\OperatorTok{);}
        \FunctionTok{loop}\OperatorTok{();} \CommentTok{//Il loop qui non riprenderà, perchè una volta che abbiamo premuto il tasto sinistro per interrompere il Loop allora usciremo anche dalla funzine \textquotesingle{}draw()\textquotesingle{}, per questo allora dovremo usare le funzioni apposite. Mostrerò il funzionamento nel codice qui sotto }
    \OperatorTok{\}}
  \OperatorTok{\}}

\OperatorTok{\}}
\end{Highlighting}
\end{Shaded}

\begin{quote}
Considerando che veniamo dal C/C++, e non siamo esperti di Java, è
importante chiarire la sintassi della dichiarazione di un array, e di
una piccolezza nella condizione del for.\\
Ecco spiegato tutto.

\begin{enumerate}
\def\labelenumi{\arabic{enumi}.}
\item
  La Dichiarazione e Allocazione
\end{enumerate}

\begin{Shaded}
\begin{Highlighting}[]
\DataTypeTok{float}\OperatorTok{[]}\NormalTok{ d }\OperatorTok{=} \KeywordTok{new} \DataTypeTok{float}\OperatorTok{[}\DecValTok{10}\OperatorTok{];}
\end{Highlighting}
\end{Shaded}

In C, per fare la stessa cosa dinamicamente, scriveresti qualcosa come:

\begin{Shaded}
\begin{Highlighting}[]
\CommentTok{// C equivalente (approssimativo)}
\DataTypeTok{float} \OperatorTok{*}\NormalTok{d }\OperatorTok{=} \OperatorTok{(}\DataTypeTok{float}\OperatorTok{*)}\NormalTok{ calloc}\OperatorTok{(}\DecValTok{10}\OperatorTok{,} \KeywordTok{sizeof}\OperatorTok{(}\DataTypeTok{float}\OperatorTok{));}
\end{Highlighting}
\end{Shaded}

Ecco le differenze fondamentali:

\begin{enumerate}
\def\labelenumi{\arabic{enumi}.}
\item
  \textbf{\texttt{float{[}{]}\ d} (Il Puntatore)}:

  \begin{itemize}
  \item
    In C scriveresti \texttt{float\ d{[}{]}} o \texttt{float\ *d}.
  \item
    In Processing, le parentesi quadre \texttt{{[}{]}} si mettono
    preferibilmente \emph{dopo il tipo} (\texttt{float{[}{]}}) per dire
    "Sto dichiarando un riferimento (puntatore) a un array di float".
  \item
    In questo momento \texttt{d} è solo un riferimento, non contiene
    ancora i dati.
  \end{itemize}
\item
  \textbf{\texttt{new} (Il Malloc Potenziato)}:

  \begin{itemize}
  \item
    Questa è la parola chiave magica. In C usi \texttt{malloc}. In
    Java/Processing usi \texttt{new}.
  \item
    \textbf{Differenza cruciale:} \texttt{new} non solo alloca la
    memoria nello \emph{Heap}, ma la \textbf{pulisce anche}.
  \item
    Mentre in C un array creato nello stack contiene valori
    "spazzatura", in Processing \texttt{new\ float{[}10{]}} crea 10
    caselle e le riempie tutte automaticamente con
    \textbf{\texttt{0.0}}.
  \end{itemize}
\item
  \textbf{\texttt{{[}10{]}} (La Dimensione)}:

  \begin{itemize}
  \item
    Stai chiedendo spazio per 10 elementi. Una volta creato,
    l\textquotesingle array \textbf{non può cambiare dimensione}
    (proprio come un blocco \texttt{malloc} in C, a meno che non usi
    \texttt{realloc}, cosa che qui non esiste direttamente).
  \end{itemize}
\end{enumerate}

\begin{center}\rule{0.5\linewidth}{0.5pt}\end{center}

\mbox{}%
\subsubsection{\texorpdfstring{2. La gestione della dimensione
(\texttt{.length})}{2. La gestione della dimensione (.length)}}\label{2-la-gestione-della-dimensione-length}

\begin{Shaded}
\begin{Highlighting}[]
\ControlFlowTok{for}\OperatorTok{(}\DataTypeTok{int}\NormalTok{ i }\OperatorTok{=} \DecValTok{0}\OperatorTok{;}\NormalTok{ i }\OperatorTok{\textless{}}\NormalTok{ d}\OperatorTok{.}\FunctionTok{length}\OperatorTok{;}\NormalTok{ i}\OperatorTok{++)} \OperatorTok{\{} \KeywordTok{...} \OperatorTok{\}}
\end{Highlighting}
\end{Shaded}

In C, un array è "stupido": non sa quanto è grande. Devi sempre portarti
dietro una variabile \texttt{int\ size} o fare calcoli con
\texttt{sizeof}.

In Processing, gli array sono \textbf{oggetti intelligenti}.

\begin{itemize}
\item
  L\textquotesingle array \texttt{d} si porta dietro "nello zaino" la
  sua dimensione.
\item
  \texttt{d.length} è una proprietà di sola lettura che ti restituisce
  \texttt{10}.
\item
  È comodissimo: non rischi mai di andare fuori dai bordi (Buffer
  Overflow) senza accorgertene, perché se provi a leggere
  \texttt{d{[}11{]}}, Processing ti ferma subito con un errore, invece
  di leggere memoria a caso come farebbe C.
\end{itemize}

\begin{center}\rule{0.5\linewidth}{0.5pt}\end{center}

\mbox{}%
\subsubsection{Riassunto C vs
Processing}\label{riassunto-c-vs-processing}

{\def\LTcaptype{none} % do not increment counter
\begin{longtable}[]{@{}lll@{}}
\toprule\noalign{}
\textbf{Concetto} & \textbf{C (classico)} & \textbf{Processing
(Java)} \\
\midrule\noalign{}
\endhead
\bottomrule\noalign{}
\endlastfoot
\textbf{Dichiarazione} & \texttt{float\ d{[}10{]};} (Stack) o
\texttt{float\ *d;} & \texttt{float{[}{]}\ d;} \\
\textbf{Allocazione} & \texttt{malloc(10\ *\ sizeof(float))} &
\texttt{new\ float{[}10{]}} \\
\textbf{Pulizia memoria} & Manuale (\texttt{memset} o \texttt{calloc}) &
Automatica (tutto a 0) \\
\textbf{Dimensione} & Devi ricordartela tu (\texttt{sizeof} o var extra)
& \texttt{d.length} (l\textquotesingle array lo sa) \\
\textbf{Deallocazione} & \texttt{free(d)} (Obbligatorio!) &
\textbf{Automatica} (Garbage Collector) 🎉 \\
\end{longtable}
}
\end{quote}

\begin{Shaded}
\begin{Highlighting}[]
\CommentTok{//[DECIMO FILE]}
\CommentTok{//creiamo un array di tipo float globalmente in modo che è accessibile ovunque}
\CommentTok{//new è l\textquotesingle{}equivalente del malloc in C, ed in questo caso abbiamo allocato un array di 10 valori, ossia 10 cerchi}
\DataTypeTok{float} \OperatorTok{[]}\NormalTok{ d }\OperatorTok{=} \KeywordTok{new} \DataTypeTok{float} \OperatorTok{[}\DecValTok{10}\OperatorTok{];} 

\DataTypeTok{void} \FunctionTok{setup}\OperatorTok{()\{}
  \FunctionTok{size}\OperatorTok{(}\DecValTok{800}\OperatorTok{,} \DecValTok{800}\OperatorTok{);}
  \ControlFlowTok{for} \OperatorTok{(}\DataTypeTok{int}\NormalTok{ i }\OperatorTok{=} \DecValTok{0}\OperatorTok{;}\NormalTok{ i }\OperatorTok{\textless{}}\NormalTok{ d}\OperatorTok{.}\FunctionTok{length}\OperatorTok{;}\NormalTok{ i}\OperatorTok{++)\{} \CommentTok{//riempiamo l\textquotesingle{}array}
\NormalTok{    d}\OperatorTok{[}\NormalTok{i}\OperatorTok{]} \OperatorTok{=} \OperatorTok{(}\NormalTok{i}\OperatorTok{+}\DecValTok{1}\OperatorTok{)*}\DecValTok{60}\OperatorTok{;} \CommentTok{//ogni valore sarà quello precedente + 60}
  \OperatorTok{\}}
\OperatorTok{\}}

\DataTypeTok{void} \FunctionTok{draw}\OperatorTok{()\{}
  \ControlFlowTok{for} \OperatorTok{(}\DataTypeTok{int}\NormalTok{ i }\OperatorTok{=}\NormalTok{ d}\OperatorTok{.}\FunctionTok{length}\OperatorTok{{-}}\DecValTok{1}\OperatorTok{;}\NormalTok{ i}\OperatorTok{\textgreater{}=}\DecValTok{0}\OperatorTok{;}\NormalTok{ i}\OperatorTok{{-}{-})\{} \CommentTok{//Scorriamo l\textquotesingle{}array al contrario altrimenti i cerchi più grandi saranno disegnati su quelli più piccoli }
    \FunctionTok{fill}\OperatorTok{(}\FunctionTok{random}\OperatorTok{(}\DecValTok{255}\OperatorTok{),} \FunctionTok{random}\OperatorTok{(}\DecValTok{255}\OperatorTok{),} \FunctionTok{random}\OperatorTok{(}\DecValTok{255}\OperatorTok{));} \CommentTok{//Riempiamo con colori random}
    \FunctionTok{ellipse}\OperatorTok{(}\NormalTok{width}\OperatorTok{/}\DecValTok{2}\OperatorTok{,}\NormalTok{ height}\OperatorTok{/}\DecValTok{2}\OperatorTok{,}\NormalTok{ d}\OperatorTok{[}\NormalTok{i}\OperatorTok{],}\NormalTok{ d}\OperatorTok{[}\NormalTok{i}\OperatorTok{]);} \CommentTok{//Il centro di ogni ellisse sarà al centro della canvas, e i diametri avranno i valori che andiamo a leggere dall\textquotesingle{}array}
  \OperatorTok{\}}
  \OperatorTok{\}}


  \CommentTok{//{-}{-}{-}FUNZIONI PER MOUSE{-}{-}{-}}
  \DataTypeTok{void} \FunctionTok{mousePressed}\OperatorTok{()\{} \CommentTok{//Adesso la nostra animazione riprenderà tranquillamente premendo tasto destro e si interromperà premendo tasto sinistro}
     \ControlFlowTok{if}\OperatorTok{(}\NormalTok{mouseButton }\OperatorTok{==}\NormalTok{ LEFT}\OperatorTok{)\{}
        \FunctionTok{println}\OperatorTok{(}\StringTok{"L"}\OperatorTok{);}
        \FunctionTok{noLoop}\OperatorTok{();}
    \OperatorTok{\}}
    \ControlFlowTok{if}\OperatorTok{(}\NormalTok{mouseButton }\OperatorTok{==}\NormalTok{ RIGHT}\OperatorTok{)\{}
        \FunctionTok{println}\OperatorTok{(}\StringTok{"R"}\OperatorTok{);}
        \FunctionTok{loop}\OperatorTok{();}
    \OperatorTok{\}}
  \OperatorTok{\}}
  \DataTypeTok{void} \FunctionTok{mouseReleased}\OperatorTok{()\{}
  \OperatorTok{\}}
  \DataTypeTok{void} \FunctionTok{mouseClicked}\OperatorTok{()\{}
  \OperatorTok{\}}

  \CommentTok{//{-}{-}{-}FUNZIONI PER TASTIERA{-}{-}{-}}
  \DataTypeTok{void} \FunctionTok{keyPressed}\OperatorTok{()\{} \CommentTok{//Vale anche se holdiamo il tasto}
  \OperatorTok{\}}
  \DataTypeTok{void} \FunctionTok{keyReleased}\OperatorTok{()\{}
  \OperatorTok{\}}
  \DataTypeTok{void} \FunctionTok{keyTyped}\OperatorTok{()\{}
  \OperatorTok{\}}
\end{Highlighting}
\end{Shaded}

\begin{Shaded}
\begin{Highlighting}[]
\CommentTok{//[UNIDCESIMO FILE]}
\CommentTok{//Adesso: Vogliamo eliminare il funzionamento di tasto destro e tasto sinistreo per interrompere e riprendere l\textquotesingle{}animazione. E vogliamo utilizzare un solo tasto come se fosse un interruttore.}
\DataTypeTok{float} \OperatorTok{[]}\NormalTok{ d }\OperatorTok{=} \KeywordTok{new} \DataTypeTok{float} \OperatorTok{[}\DecValTok{10}\OperatorTok{];} 
\CommentTok{//utilizziamo una variabile booleana}
\DataTypeTok{boolean}\NormalTok{ vai }\OperatorTok{=} \KeywordTok{true}\OperatorTok{;}

\DataTypeTok{void} \FunctionTok{setup}\OperatorTok{()\{}
  \FunctionTok{size}\OperatorTok{(}\DecValTok{800}\OperatorTok{,} \DecValTok{800}\OperatorTok{);}
  \ControlFlowTok{for} \OperatorTok{(}\DataTypeTok{int}\NormalTok{ i }\OperatorTok{=} \DecValTok{0}\OperatorTok{;}\NormalTok{ i }\OperatorTok{\textless{}}\NormalTok{ d}\OperatorTok{.}\FunctionTok{length}\OperatorTok{;}\NormalTok{ i}\OperatorTok{++)\{}
\NormalTok{    d}\OperatorTok{[}\NormalTok{i}\OperatorTok{]} \OperatorTok{=} \OperatorTok{(}\NormalTok{i}\OperatorTok{+}\DecValTok{1}\OperatorTok{)*}\DecValTok{60}\OperatorTok{;}
  \OperatorTok{\}}
\OperatorTok{\}}

\DataTypeTok{void} \FunctionTok{draw}\OperatorTok{()\{}
  \ControlFlowTok{for} \OperatorTok{(}\DataTypeTok{int}\NormalTok{ i }\OperatorTok{=}\NormalTok{ d}\OperatorTok{.}\FunctionTok{length}\OperatorTok{{-}}\DecValTok{1}\OperatorTok{;}\NormalTok{ i}\OperatorTok{\textgreater{}=}\DecValTok{0}\OperatorTok{;}\NormalTok{ i}\OperatorTok{{-}{-})\{}
    \FunctionTok{fill}\OperatorTok{(}\FunctionTok{random}\OperatorTok{(}\DecValTok{255}\OperatorTok{),} \FunctionTok{random}\OperatorTok{(}\DecValTok{255}\OperatorTok{),} \FunctionTok{random}\OperatorTok{(}\DecValTok{255}\OperatorTok{));}
    \FunctionTok{ellipse}\OperatorTok{(}\NormalTok{width}\OperatorTok{/}\DecValTok{2}\OperatorTok{,}\NormalTok{ height}\OperatorTok{/}\DecValTok{2}\OperatorTok{,}\NormalTok{ d}\OperatorTok{[}\NormalTok{i}\OperatorTok{],}\NormalTok{ d}\OperatorTok{[}\NormalTok{i}\OperatorTok{]);}
  \OperatorTok{\}}
  \OperatorTok{\}}

  \DataTypeTok{void} \FunctionTok{mousePressed}\OperatorTok{()\{}
  \OperatorTok{\}}
  \DataTypeTok{void} \FunctionTok{mouseReleased}\OperatorTok{()\{}
  \OperatorTok{\}}
  \DataTypeTok{void} \FunctionTok{mouseClicked}\OperatorTok{()\{}
    \CommentTok{//Qui stiamo dicendo che la nostra animazione si attacca e si stacca, con qualunque tasto del mouse (sinistro, centro, destro)}
\NormalTok{    vai }\OperatorTok{=} \OperatorTok{!}\NormalTok{vai}\OperatorTok{;}
    \ControlFlowTok{if}\OperatorTok{(}\NormalTok{vai }\OperatorTok{==} \KeywordTok{true}\OperatorTok{)\{}
        \FunctionTok{loop}\OperatorTok{();}
    \OperatorTok{\}}\ControlFlowTok{else}\OperatorTok{\{}
        \FunctionTok{noLoop}\OperatorTok{();}
    \OperatorTok{\}}
  \OperatorTok{\}}

  \DataTypeTok{void} \FunctionTok{keyPressed}\OperatorTok{()\{}
  \OperatorTok{\}}
  \DataTypeTok{void} \FunctionTok{keyReleased}\OperatorTok{()\{}
  \OperatorTok{\}}
  \DataTypeTok{void} \FunctionTok{keyTyped}\OperatorTok{()\{}
  \OperatorTok{\}}
\end{Highlighting}
\end{Shaded}

\begin{Shaded}
\begin{Highlighting}[]
\CommentTok{//[DODICESIMO FILE]}
\CommentTok{//Adesso lavoriamo sull\textquotesingle{}utilizzo delle freccette della tastiera}
\DataTypeTok{int}\NormalTok{ x }\OperatorTok{=} \DecValTok{0}\OperatorTok{;}

\DataTypeTok{void} \FunctionTok{setup}\OperatorTok{()\{}
   \FunctionTok{size}\OperatorTok{(}\DecValTok{500}\OperatorTok{,}\DecValTok{500}\OperatorTok{);} 
\OperatorTok{\}}

\DataTypeTok{void} \FunctionTok{draw}\OperatorTok{()\{}
    \FunctionTok{background}\OperatorTok{(}\DecValTok{0}\OperatorTok{);}
    \FunctionTok{rect}\OperatorTok{(}\NormalTok{x}\OperatorTok{,}\NormalTok{height}\OperatorTok{/}\DecValTok{2}\OperatorTok{,} \DecValTok{50}\OperatorTok{,}\DecValTok{50}\OperatorTok{);}
    \ControlFlowTok{if}\OperatorTok{(}\NormalTok{keyPressed}\OperatorTok{)\{}
        \ControlFlowTok{if}\OperatorTok{(}\NormalTok{key}\OperatorTok{==}\NormalTok{CODED}\OperatorTok{)\{}
        \ControlFlowTok{if}\OperatorTok{(}\NormalTok{keyCode}\OperatorTok{==}\NormalTok{LEFT}\OperatorTok{)\{} \CommentTok{//Se il keyCode corrisponde al tasto della freccia sinistra, decrementiamo il valore di x.}
\NormalTok{            x}\OperatorTok{{-}=}\DecValTok{5}\OperatorTok{;}
        \OperatorTok{\}}
        \ControlFlowTok{if}\OperatorTok{(}\NormalTok{keyCode}\OperatorTok{==}\NormalTok{RIGHT}\OperatorTok{)\{} \CommentTok{//Se il keyCode corrisponde al tasto della freccia destra, incrementiamo il valore di x.}
\NormalTok{            x}\OperatorTok{+=}\DecValTok{5}\OperatorTok{;}
        \OperatorTok{\}}
    \OperatorTok{\}}
    \OperatorTok{\}}
\OperatorTok{\}}

\DataTypeTok{void} \FunctionTok{keyPressed}\OperatorTok{()\{} \CommentTok{//Ci sono dei caratteri che sono codificati, detti \textquotesingle{}CODED\textquotesingle{}.}
\CommentTok{/*    if(key==CODED)\{}
\CommentTok{        if(keyCode==LEFT)\{ //Se il keyCode corrisponde al tasto della freccia sinistra, decrementiamo il valore di x.}
\CommentTok{            x{-}=5;}
\CommentTok{        \}}
\CommentTok{        if(keyCode==RIGHT)\{ //Se il keyCode corrisponde al tasto della freccia destra, incrementiamo il valore di x.}
\CommentTok{            x+=5;}
\CommentTok{        \}}
\CommentTok{    \}*/}
\OperatorTok{\}} 

\CommentTok{//HO CONSTATATO CHE A QUANTO PARE USARLO NEL DRAW è MEGLIO. OVVIAMENTE BISOGNA USARE IL BOOLEANO \textquotesingle{}keyPressed\textquotesingle{} PER S}
\CommentTok{/*}
\CommentTok{Codici dei tasti CODED reperibili con keyCode;}
\CommentTok{UP, DOWN, LEFT, RIGHT, ALT, CONTROL e SHIFT}

\CommentTok{Mentre in key, non avendo bisogno di fare il controlo se key==CODED, abbiamo i seguenti caratteri speciali:}
\CommentTok{BACKSPACE, TAB, ENTER, RETURN, ESC, e DELETE.}
\CommentTok{L\textquotesingle{}invio si può "catturare" in due modi (key == ENTER) oppure (key == \textquotesingle{}\textbackslash{}n\textquotesingle{}).}
\CommentTok{Lo spazio è il carattere \textquotesingle{} \textquotesingle{} (c\textquotesingle{}è uno spazio tra i singoli apici)}
\CommentTok{*/}
\end{Highlighting}
\end{Shaded}

\begin{Shaded}
\begin{Highlighting}[]
\CommentTok{//[TREDICESIMO FILE]}
\CommentTok{/*}
\CommentTok{    Esercizio di PacMan animato. Il personaggio dovrà aprire e chiudere la bocca e tornare avanti ed indietro.}
\CommentTok{*/}

\CommentTok{//Dichiariamo le variabili che ci serviranno}
\DataTypeTok{int}\NormalTok{ x }\OperatorTok{=} \DecValTok{60}\OperatorTok{;}
\DataTypeTok{int}\NormalTok{ dir }\OperatorTok{=} \DecValTok{1}\OperatorTok{;} \CommentTok{//direzione}
\DataTypeTok{int}\NormalTok{ a }\OperatorTok{=} \DecValTok{30}\OperatorTok{;} \CommentTok{//angolo = 30°}
\DataTypeTok{int}\NormalTok{ dira }\OperatorTok{=} \OperatorTok{{-}}\DecValTok{1}\OperatorTok{;}
\DataTypeTok{int}\NormalTok{ v }\OperatorTok{=} \DecValTok{10}\OperatorTok{;} \CommentTok{//velocità}
\DataTypeTok{int}\NormalTok{ va }\OperatorTok{=}\NormalTok{ v}\OperatorTok{/}\DecValTok{2}\OperatorTok{;} \CommentTok{//velocità di apertura bocca}
\DataTypeTok{int}\NormalTok{ r }\OperatorTok{=} \DecValTok{40}\OperatorTok{;} \CommentTok{//raggio}

\DataTypeTok{void} \FunctionTok{setup}\OperatorTok{()\{}
    \FunctionTok{size}\OperatorTok{(}\DecValTok{500}\OperatorTok{,}\DecValTok{500}\OperatorTok{);}
    \FunctionTok{ellipseMode}\OperatorTok{(}\NormalTok{RADIUS}\OperatorTok{);}  \CommentTok{// CAMBIA IL MODO DI DISEGNARE: I parametri diventano il RAGGIO (metà larghezza), non il DIAMETRO. Senza questa riga, il Pacman sarebbe largo 40 totali (piccolo). Con questa riga, 40 è solo la metà, quindi il totale è 80 (grande il doppio).}
\OperatorTok{\}}

\DataTypeTok{void} \FunctionTok{draw}\OperatorTok{()\{}
    \FunctionTok{background}\OperatorTok{(}\DecValTok{0}\OperatorTok{);} \CommentTok{//Impostiamo lo sfondo nero}
    \FunctionTok{fill}\OperatorTok{(}\DecValTok{255}\OperatorTok{,} \DecValTok{255}\OperatorTok{,} \DecValTok{0}\OperatorTok{);} \CommentTok{//Coloriamo di giallo il nostro pacman}
    \ControlFlowTok{if}\OperatorTok{(}\NormalTok{dir}\OperatorTok{==}\DecValTok{1}\OperatorTok{)\{} \CommentTok{//Se il nostro pacman si muove verso destra...}
    \FunctionTok{arc}\OperatorTok{(}\NormalTok{x}\OperatorTok{,}\NormalTok{height}\OperatorTok{/}\DecValTok{2}\OperatorTok{,}\NormalTok{ r}\OperatorTok{,}\NormalTok{r}\OperatorTok{,} \FunctionTok{radians}\OperatorTok{(}\NormalTok{a}\OperatorTok{),} \FunctionTok{radians}\OperatorTok{(}\DecValTok{360}\OperatorTok{{-}}\NormalTok{a}\OperatorTok{));} \CommentTok{//...allora la bocca sarà a destra}
    \OperatorTok{\}}
    \ControlFlowTok{if}\OperatorTok{(}\NormalTok{dir}\OperatorTok{=={-}}\DecValTok{1}\OperatorTok{)\{} \CommentTok{//Se il nostro pacman si muove verso destra...}
        \FunctionTok{arc}\OperatorTok{(}\NormalTok{x}\OperatorTok{,}\NormalTok{height}\OperatorTok{/}\DecValTok{2}\OperatorTok{,}\NormalTok{ r}\OperatorTok{,}\NormalTok{r}\OperatorTok{,} \FunctionTok{radians}\OperatorTok{(}\DecValTok{180}\OperatorTok{+}\NormalTok{a}\OperatorTok{),} \FunctionTok{radians}\OperatorTok{(}\DecValTok{360}\OperatorTok{+}\DecValTok{180}\OperatorTok{{-}}\NormalTok{a}\OperatorTok{));} \CommentTok{//...allora la bocca sarà a sinistra}
    \OperatorTok{\}}
    \FunctionTok{fill}\OperatorTok{(}\DecValTok{0}\OperatorTok{);} \CommentTok{//Colore dell\textquotesingle{}occhio}
    \FunctionTok{ellipse}\OperatorTok{(}\NormalTok{x}\OperatorTok{,}\NormalTok{ height}\OperatorTok{/}\DecValTok{2}\OperatorTok{{-}}\DecValTok{20}\OperatorTok{,} \DecValTok{5}\OperatorTok{,}\DecValTok{5}\OperatorTok{);} \CommentTok{//Definizione dell\textquotesingle{}ellisse che sarà l\textquotesingle{}occhio}
    \CommentTok{//Questo è il controllo dell\textquotesingle{}occhio. }
    \CommentTok{//Se x è maggiore uguale della larghezza della canvas meno il raggio (si aggiunge {-} r perché dobbiamo controllare che sia la bocca a toccare e non il centro)...7}
    \CommentTok{//Oppure se x è minore del raggio}
    \CommentTok{//allora cambia la direzione}
    \ControlFlowTok{if}\OperatorTok{(}\NormalTok{x }\OperatorTok{\textgreater{}=}\NormalTok{ width}\OperatorTok{{-}}\NormalTok{r }\OperatorTok{||}\NormalTok{ x }\OperatorTok{\textless{}}\NormalTok{ r}\OperatorTok{)\{} 
\NormalTok{            dir }\OperatorTok{=} \OperatorTok{{-}}\NormalTok{dir}\OperatorTok{;} \CommentTok{//dir = {-}dir è un interruttore. È un modo velocissimo per dire al computer: "Qualunque cosa tu sia (1 o {-}1), diventa il tuo opposto".}
    \OperatorTok{\}}
    \ControlFlowTok{if}\OperatorTok{(}\NormalTok{a }\OperatorTok{\textgreater{}} \DecValTok{30} \OperatorTok{||}\NormalTok{ a }\OperatorTok{\textless{}=} \DecValTok{0}\OperatorTok{)\{} \CommentTok{//Questo comanda l\textquotesingle{}apertura della bocca.}
    \CommentTok{//Se a \textgreater{} 30 la bocca è aperta}
    \CommentTok{//Se a \textless{}= 0 la bocca è chiusa}
\NormalTok{        dira }\OperatorTok{=} \OperatorTok{{-}}\NormalTok{dira}\OperatorTok{;}
        \CommentTok{/*Immagina dira come un interruttore del respiro.}
\CommentTok{        Quando dira è positivo (es. +1): Pacman sta INSPIRANDO. Ad ogni fotogramma aggiungi 1 all\textquotesingle{}angolo. La bocca si apre sempre di più. Angolo: 10°... 20°... 30° (Bocca spalancata! Aaaah!)}
\CommentTok{        Quando dira è negativo (es. {-}1): Pacman sta ESPIRANDO. Ad ogni fotogramma togli 1 all\textquotesingle{}angolo. La bocca si chiude sempre di più. Angolo: 30°... 20°... 10°... 0° (Bocca chiusa! Mmh!)}
\CommentTok{        Quindi cosa fa dira = {-}dira? È il comando che dice: "Inverti il respiro!"}
\CommentTok{        Il computer controlla continuamente:}
\CommentTok{        "Ehi, la bocca è arrivata a 30 (troppo aperta)?" {-}\textgreater{} SÌ! {-}\textgreater{} Allora fai dira = {-}dira. Da ora in poi chiudi.}
\CommentTok{        (Il tempo passa, la bocca si chiude...)}
\CommentTok{        "Ehi, la bocca è arrivata a 0 (troppo chiusa)?" {-}\textgreater{} SÌ! {-}\textgreater{} Allora fai dira = {-}dira. Da ora in poi apri.}
\CommentTok{        È un ciclo infinito: Apri {-}\textgreater{} Inverti {-}\textgreater{} Chiudi {-}\textgreater{} Inverti {-}\textgreater{} Apri... Proprio come un polmone o un cuore che batte.*/}
    \OperatorTok{\}}
\NormalTok{    a}\OperatorTok{+=}\NormalTok{dira}\OperatorTok{*}\NormalTok{va}\OperatorTok{;} \CommentTok{//Ad ogni fotogrammala bocca si muove un po\textquotesingle{} (e con va comandiamo la velocità della bocca, che è proporzionale a v)}
\NormalTok{    x }\OperatorTok{+=}\NormalTok{ dir }\OperatorTok{*}\NormalTok{ v}\OperatorTok{;} \CommentTok{//Prendi la posizione x e aggiungi la velocità v nella direzione giusta dir. Questa riga cambia la direzione, la velocità di movimento è sempre v}
\OperatorTok{\}}
\end{Highlighting}
\end{Shaded}

\begin{Shaded}
\begin{Highlighting}[]
\CommentTok{//[QUATTORDICESIMO FILE]}
\CommentTok{//Introduciamo l\textquotesingle{}interpolazione in processing: \textquotesingle{}lerp()\textquotesingle{} [Linear Interpolation]}
\CommentTok{//Gli forniamo dei key frame, lui si muoverà da un punto ad un altro calcolando dei punti intermedi, dunque ci servirà per l\textquotesingle{}animazione.}
\CommentTok{//Come parametri prende: un punto di inizio, un punto di fine e la percentuale. Per un totale di 3 parametri.}

\CommentTok{//Proviamo prima su un disegno e poi su una animazione.}

\DataTypeTok{int}\NormalTok{ x1 }\OperatorTok{=} \DecValTok{15}\OperatorTok{;}
\DataTypeTok{int}\NormalTok{ y1 }\OperatorTok{=} \DecValTok{10}\OperatorTok{;}
\DataTypeTok{int}\NormalTok{ x2 }\OperatorTok{=} \DecValTok{80}\OperatorTok{;}
\DataTypeTok{int}\NormalTok{ y2 }\OperatorTok{=} \DecValTok{90}\OperatorTok{;}
\FunctionTok{line}\OperatorTok{(}\NormalTok{x1}\OperatorTok{,}\NormalTok{y1}\OperatorTok{,}\NormalTok{x2}\OperatorTok{,}\NormalTok{y2}\OperatorTok{);}
\ControlFlowTok{for}\OperatorTok{(}\DataTypeTok{int}\NormalTok{ i }\OperatorTok{=} \DecValTok{0}\OperatorTok{;}\NormalTok{ i }\OperatorTok{\textless{}=} \DecValTok{10}\OperatorTok{;}\NormalTok{ i}\OperatorTok{++)\{} 
    \FunctionTok{stroke}\OperatorTok{(}\DecValTok{255}\OperatorTok{,} \DecValTok{0}\OperatorTok{,} \DecValTok{0}\OperatorTok{);}
    \FunctionTok{strokeWeight}\OperatorTok{(}\DecValTok{3}\OperatorTok{);}
    \DataTypeTok{float}\NormalTok{ x }\OperatorTok{=} \FunctionTok{lerp}\OperatorTok{(}\NormalTok{x1}\OperatorTok{,}\NormalTok{x2}\OperatorTok{,}\NormalTok{ i}\OperatorTok{/}\FloatTok{10.0}\OperatorTok{);} 
    \DataTypeTok{float}\NormalTok{ y }\OperatorTok{=} \FunctionTok{lerp}\OperatorTok{(}\NormalTok{y1}\OperatorTok{,}\NormalTok{y2}\OperatorTok{,}\NormalTok{ i}\OperatorTok{/}\FloatTok{10.0}\OperatorTok{);}
    \FunctionTok{point}\OperatorTok{(}\NormalTok{x}\OperatorTok{,}\NormalTok{y}\OperatorTok{);}
\OperatorTok{\}}

\CommentTok{//Quindi abbiamo disegnato dei puntini al 10\%, 20\%, 30\%, 40\%, 50\%, 60\%, 70\%, 80\%, 90\%, 100\%.}
\end{Highlighting}
\end{Shaded}

\begin{Shaded}
\begin{Highlighting}[]
\CommentTok{//[QUINDICESIMO FILE]}
 \CommentTok{//Adesso utilizziamo la funzione di interpolazione nelle animazioni}
\DataTypeTok{float}\NormalTok{ x}\OperatorTok{;}
\DataTypeTok{float}\NormalTok{ y}\OperatorTok{;}

\DataTypeTok{void} \FunctionTok{setup}\OperatorTok{()\{}
    \FunctionTok{size}\OperatorTok{(}\DecValTok{500}\OperatorTok{,}\DecValTok{500}\OperatorTok{);}
    \FunctionTok{noStroke}\OperatorTok{();}
\OperatorTok{\}}

\DataTypeTok{void} \FunctionTok{draw}\OperatorTok{()\{}
    \FunctionTok{background}\OperatorTok{(}\DecValTok{0}\OperatorTok{);}
    \CommentTok{//x=mouseX;  }
    \CommentTok{//y=mouseY;}
\NormalTok{    x }\OperatorTok{=} \FunctionTok{lerp}\OperatorTok{(}\NormalTok{x}\OperatorTok{,}\NormalTok{ mouseX}\OperatorTok{,} \FloatTok{0.1}\OperatorTok{);} \CommentTok{//0.1 riguarda lo spostamento}
\NormalTok{    y }\OperatorTok{=} \FunctionTok{lerp}\OperatorTok{(}\NormalTok{y}\OperatorTok{,}\NormalTok{ mouseY}\OperatorTok{,} \FloatTok{0.1}\OperatorTok{);}
    \FunctionTok{ellipse}\OperatorTok{(}\NormalTok{x}\OperatorTok{,}\NormalTok{y}\OperatorTok{,}\DecValTok{60}\OperatorTok{,}\DecValTok{60}\OperatorTok{);}
\OperatorTok{\}}
\end{Highlighting}
\end{Shaded}

\begin{Shaded}
\begin{Highlighting}[]
\CommentTok{//[SEDICESIMO FILE]}
\CommentTok{//Facciamo una pallina che lascia la scia}

\DataTypeTok{float}\NormalTok{ x}\OperatorTok{;}
\DataTypeTok{float}\NormalTok{ y}\OperatorTok{;}

\DataTypeTok{void} \FunctionTok{setup}\OperatorTok{()\{}
    \FunctionTok{size}\OperatorTok{(}\DecValTok{500}\OperatorTok{,}\DecValTok{500}\OperatorTok{);}
    \FunctionTok{noStroke}\OperatorTok{();}
\OperatorTok{\}}

\DataTypeTok{void} \FunctionTok{draw}\OperatorTok{()\{}
    \CommentTok{//background(0);}
    \FunctionTok{fill}\OperatorTok{(}\DecValTok{0}\OperatorTok{,} \DecValTok{40}\OperatorTok{);}
    \FunctionTok{rect}\OperatorTok{(}\DecValTok{0}\OperatorTok{,} \DecValTok{0}\OperatorTok{,}\NormalTok{ width}\OperatorTok{,}\NormalTok{ height}\OperatorTok{);}
\NormalTok{    x }\OperatorTok{=} \FunctionTok{lerp}\OperatorTok{(}\NormalTok{x}\OperatorTok{,}\NormalTok{ mouseX}\OperatorTok{,} \FloatTok{0.1}\OperatorTok{);}
\NormalTok{    y }\OperatorTok{=} \FunctionTok{lerp}\OperatorTok{(}\NormalTok{y}\OperatorTok{,}\NormalTok{ mouseY}\OperatorTok{,} \FloatTok{0.1}\OperatorTok{);}
    \FunctionTok{fill}\OperatorTok{(}\DecValTok{255}\OperatorTok{);}
    \FunctionTok{ellipse}\OperatorTok{(}\NormalTok{x}\OperatorTok{,}\NormalTok{y}\OperatorTok{,}\DecValTok{60}\OperatorTok{,}\DecValTok{60}\OperatorTok{);}
\OperatorTok{\}}
\end{Highlighting}
\end{Shaded}

\begin{Shaded}
\begin{Highlighting}[]
\CommentTok{//[DICIASETTESIMO FILE]}
 \CommentTok{//Movimento randomico di una pallina}

\DataTypeTok{float}\NormalTok{ x}\OperatorTok{;}
\DataTypeTok{float}\NormalTok{ y}\OperatorTok{;}

\DataTypeTok{void} \FunctionTok{setup}\OperatorTok{()\{}
    \FunctionTok{size}\OperatorTok{(}\DecValTok{500}\OperatorTok{,}\DecValTok{500}\OperatorTok{);}

\NormalTok{    x}\OperatorTok{=}\NormalTok{width}\OperatorTok{/}\DecValTok{2}\OperatorTok{;}
\NormalTok{    y}\OperatorTok{=}\NormalTok{height}\OperatorTok{/}\DecValTok{2}\OperatorTok{;}
    \FunctionTok{noStroke}\OperatorTok{();}
\OperatorTok{\}}

\DataTypeTok{void} \FunctionTok{draw}\OperatorTok{()\{}
    \CommentTok{//background(0);}
    \FunctionTok{fill}\OperatorTok{(}\DecValTok{0}\OperatorTok{,} \DecValTok{25}\OperatorTok{);}
    \FunctionTok{rect}\OperatorTok{(}\DecValTok{0}\OperatorTok{,} \DecValTok{0}\OperatorTok{,}\NormalTok{ width}\OperatorTok{,}\NormalTok{ height}\OperatorTok{);}
\NormalTok{    x }\OperatorTok{=}\NormalTok{ x}\OperatorTok{+}\FunctionTok{random}\OperatorTok{({-}}\DecValTok{5}\OperatorTok{,}\DecValTok{5}\OperatorTok{);}
\NormalTok{    y }\OperatorTok{=}\NormalTok{ y}\OperatorTok{+}\FunctionTok{random}\OperatorTok{({-}}\DecValTok{5}\OperatorTok{,}\DecValTok{5}\OperatorTok{);}
    \FunctionTok{fill}\OperatorTok{(}\DecValTok{255}\OperatorTok{);}
    \FunctionTok{ellipse}\OperatorTok{(}\NormalTok{x}\OperatorTok{,}\NormalTok{y}\OperatorTok{,}\DecValTok{60}\OperatorTok{,}\DecValTok{60}\OperatorTok{);}
\OperatorTok{\}}
\end{Highlighting}
\end{Shaded}

\begin{Shaded}
\begin{Highlighting}[]
\CommentTok{//[DICIOTTESIMO FILE]
}
\CommentTok{//Programmiamo una sorta di paint con 3 colori
}


\DataTypeTok{void} \FunctionTok{setup}\OperatorTok{()\{}

    \FunctionTok{size}\OperatorTok{(}\DecValTok{500}\OperatorTok{,}\DecValTok{500}\OperatorTok{);}

     \FunctionTok{stroke}\OperatorTok{(}\DecValTok{255}\OperatorTok{,} \DecValTok{0}\OperatorTok{,} \DecValTok{0}\OperatorTok{);}

    \FunctionTok{stroke}\OperatorTok{(}\DecValTok{10}\OperatorTok{);}

    \FunctionTok{background}\OperatorTok{(}\DecValTok{255}\OperatorTok{);}

\OperatorTok{\}}



\DataTypeTok{void} \FunctionTok{draw}\OperatorTok{()\{}

    \FunctionTok{strokeWeight}\OperatorTok{(}\DecValTok{5}\OperatorTok{);}

    \ControlFlowTok{if}\OperatorTok{(}\NormalTok{mousePressed}\OperatorTok{)\{}

        \CommentTok{//pmouse è una variabile che salva la posizione precedente del click del mouse
}
        \FunctionTok{line}\OperatorTok{(}\NormalTok{pmouseX}\OperatorTok{,}\NormalTok{ pmouseY}\OperatorTok{,}\NormalTok{ mouseX}\OperatorTok{,}\NormalTok{ mouseY}\OperatorTok{);}

    \OperatorTok{\}}

\OperatorTok{\}}



\DataTypeTok{void} \FunctionTok{keyPressed}\OperatorTok{()\{}

    \ControlFlowTok{if}\OperatorTok{(}\NormalTok{key }\OperatorTok{==} \CharTok{\textquotesingle{}r\textquotesingle{}} \OperatorTok{||}\NormalTok{ key }\OperatorTok{==} \CharTok{\textquotesingle{}R\textquotesingle{}}\OperatorTok{)\{}

        \FunctionTok{stroke}\OperatorTok{(}\DecValTok{255}\OperatorTok{,}\DecValTok{0}\OperatorTok{,}\DecValTok{0}\OperatorTok{);}

    \OperatorTok{\}}

    \ControlFlowTok{if}\OperatorTok{(}\NormalTok{key }\OperatorTok{==} \CharTok{\textquotesingle{}g\textquotesingle{}} \OperatorTok{||}\NormalTok{ key }\OperatorTok{==} \CharTok{\textquotesingle{}G\textquotesingle{}}\OperatorTok{)\{}

        \FunctionTok{stroke}\OperatorTok{(}\DecValTok{0}\OperatorTok{,}\DecValTok{255}\OperatorTok{,}\DecValTok{0}\OperatorTok{);}

    \OperatorTok{\}}

    \ControlFlowTok{if}\OperatorTok{(}\NormalTok{key }\OperatorTok{==} \CharTok{\textquotesingle{}b\textquotesingle{}} \OperatorTok{||}\NormalTok{ key }\OperatorTok{==} \CharTok{\textquotesingle{}B\textquotesingle{}}\OperatorTok{)\{}

        \FunctionTok{stroke}\OperatorTok{(}\DecValTok{0}\OperatorTok{,}\DecValTok{0}\OperatorTok{,}\DecValTok{255}\OperatorTok{);}

    \OperatorTok{\}}

    \ControlFlowTok{if}\OperatorTok{(}\NormalTok{key }\OperatorTok{==} \CharTok{\textquotesingle{} \textquotesingle{}}\OperatorTok{)\{}

        \FunctionTok{background}\OperatorTok{(}\DecValTok{255}\OperatorTok{);}

    \OperatorTok{\}}

\OperatorTok{\}}
\end{Highlighting}
\end{Shaded}


\end{document}
